\scanpage{11}\mspage{8}

Prouvons (iv). Supposons prouvé
\[
 \pi_{j\alpha}\in \unsure{E|\mathcal E}\qquad\text{pour }j\notin[i,2n-i]
 \tag*{(6.7)}
\]
où $i\le n$, et prouvons que
\[
 \deleted{\pi_{j\alpha}\in E\quad\text{pour } j\in}
\]
\[
 \pi_{j\alpha}\in E\quad\deleted{\text{pour}}\ \text{si }j=i\ \text{ou }j=2n-i,
 \tag*{$(6.7')$}
\]
ce qui prouvera (iv) \inserted{par récurrence ascendante sur $i\le n$.} Pour prouver 6.8. on procède par récurrence \deleted{\unsure{sur}} descendante sur $\alpha$, en supposant qu'on a
\[
 \deleted{\unsure{\pi_{j\alpha}}}\quad \pi_{i,\alpha'}\in E,\quad\pi_{2n-i,\alpha'}\in E\quad\text{si }\alpha'>\alpha,
 \tag*{(6.8)}
\]
et on en \unsure{conclura|conclure} que
\[
 \pi_{i,\alpha}\in E,\quad\pi_{2n-i,\alpha}\in E .
 \tag*{$(6.8')$}
\]
\deleted{Cela \deleted{\unsure{étant}} prouvera bien (6.7'), puisque on peut démarrer la récurrence ($\pi_{i,\alpha}=\pi_{2n-i,\alpha}=0$ si $\alpha$ grand). Si $\alpha=0$, on gagne grâce à la relation
\[
 \sum_\alpha\pi_{i,\alpha}=\pi_i\quad\text{donc}\quad\pi_{i0}=\pi_i-\sum_{\alpha\ge1}\pi_{i,\alpha},
 \tag*{(6.9)}
\]
donc il suffit de regarder le cas $\alpha\ge1$.}

Posons
\[
 q_{i,\alpha}=1-\sum_{\alpha'>\alpha}\pi_{i,\alpha'}=\sum_{\alpha'\le\alpha}\pi_{i,\alpha}
\]
donc grâce à (6.8) on a
\[
 q_{i,\alpha}\in E\ \deleted{\unsure{E}}
\]
% archive sheet number: 10
